\subsection{Topological normalization and semilocal models}

\begin{lemma}[Known; finite topological normalization, with the semilocal derivation]
\label{A:semilocal-topnormalization}
Let $R$ be a semilocal noetherian domain with fraction field $K$,
and let $\overline R$ be its integral closure in $K$.
There is a finite birational $R$-subalgebra
$R\subset C\subset\overline R$ such that
\[
 \Spec\overline R\longrightarrow\Spec C
\]
is a universal homeomorphism. Consequently there is $0\neq a\in R$
such that the normalization morphism of $\Spec R[1/a]$ is a
universal homeomorphism. Neither normalization is assumed finite
or noetherian.
\end{lemma}
\begin{proof}
Let $\mathfrak m_1,\ldots,\mathfrak m_r$ be the maximal ideals of
$R$ and set
\[
 T=\prod_{i=1}^r\widehat{R_{\mathfrak m_i}},
\]
where each completion is along its maximal ideal. The map $R\to T$
is flat, since localization and noetherian completion are flat and
the product is finite. It is faithfully flat: for a proper ideal
$I\subset R$, choose $\mathfrak m_i\supset I$; the corresponding
factor of $T/IT$ has a nonzero residue field. The ring $T$ is
noetherian. Each of its minimal-prime quotients is a complete
noetherian local domain, and hence is Japanese by
\citet[Chapitre 0, Th\'eor\`eme 23.1.5, pp.~215--216]{A_EGAIV1}.

Apply \citet[Chapitre 0, Proposition 23.2.4(ii), p.~218]{A_EGAIV1}
to the faithfully flat map $R\to T$ and the trivial finite
extension $K/K$. Its directed family consists of the finite
$R$-subalgebras of $\overline R$; each already has fraction field
$K$ because it contains $R$. The proposition gives such an
algebra $C$ for which $\Spec\overline R\to\Spec C$ is radicial.
This map is also integral and surjective, so it is a universal
homeomorphism. For a local ring this is exactly the consequence
stated in Corollaire 23.2.5, pp.~218--219, of the same source.

Choose finite $R$-algebra generators $c_1,\ldots,c_s$ of $C$ and
write them as fractions in $K=\Frac R$. The product $a\neq0$ of
their denominators gives $C[1/a]=R[1/a]$. Integral closure
commutes with localization, and universal homeomorphisms are
preserved by base change. Therefore
\[
 \Spec\overline R[1/a]\longrightarrow\Spec R[1/a]
\]
is the asserted normalization and universal homeomorphism.
\end{proof}

\begin{lemma}[Proved; the strictly local approximation bound for every overring]
\label{A:strict-overring-bound}
Let $B$ be a quasi-excellent strictly henselian noetherian local
domain, let $p$ be an invertible prime, and write
$K=\Frac B$ and $d=\dim B$. Every intermediate ring
$B\subset T\subset K$ satisfies
\[
 \cd_p(\Spec T)\leq d.
\]
The ring $T$ need not be noetherian or integrally closed.
\end{lemma}
\begin{proof}
Write $T$ as the filtered union of its finite-type $B$-subalgebras
$T_i$. Each $T_i$ is a domain containing $B$ and contained in $K$,
so $\Frac T_i=K$. Thus $\Spec T_i\to\Spec B$ is affine,
finite type and dominant with generic relative dimension zero.
The strictly local affine theorem, in the form proved and checked
in Theorem~\ref{A:relative-affine}, gives
$\cd_p(\Spec T_i)\leq d$.

The inverse system of affine schemes $\Spec T_i$ has affine
transition maps and limit $\Spec T$. A constructible torsion
\'etale sheaf descends to a finite stage, and its cohomology is
the filtered colimit of the cohomology at subsequent stages.
Consequently it has zero cohomology in degrees greater than $d$.
The coherent-site constructible approximation and coefficient
d\'evissage used in Lemma~\ref{A:A2-approximation} then give the
same vanishing for every $p$-primary torsion sheaf on $\Spec T$.
This is precisely the argument of
Proposition~\ref{A:strict-model-rigidity}, with its unnecessary
noetherianity and normality hypotheses on the intermediate ring
removed; the base $B$ still has all the hypotheses required by
the affine theorem.
\end{proof}

\begin{theorem}[Proved under stated hypotheses; A.2 for topologically normal and semilocal models]
\label{A:A2-topnormal-models}
Let $B,p,K,d$ be as in Lemma~\ref{A:strict-overring-bound}.
Let $R\subset K$ be a noetherian domain with $\Frac R=K$ and
$p\in R^\times$. Suppose that $\Spec R$ has a nonempty
geometrically unibranch open subscheme. Then A.2 holds for
$\Spec R$. In particular, A.2 holds for every semilocal
noetherian model $R$ of this field, without a normality assumption
and in every dimension.
\end{theorem}
\begin{proof}
Theorem~\ref{A:strict-field} gives $\cd_p(K)=d$.
Choose a nonzero $a\in R$ such that $\Spec R[1/a]$ lies in the
given geometrically unibranch open. For an irreducible
geometrically unibranch scheme the normalization is a universal
homeomorphism, without any finiteness hypothesis on the
normalization; this is
\citet[Tag 0GIQ, Lemma 29.55.12]{A_Stacks}.

Let $W\subset\Spec R[1/a]$ be any nonempty affine open and put
$S=\Gamma(W,\mathcal O_W)\subset K$. Its normalization
$S^\nu\subset K$ is the normalization restricted to $W$, since
normalization commutes with restriction to open subschemes.
Consequently $\Spec S^\nu\to\Spec S$ is a universal
homeomorphism. The Mori--Nagata containment of
Lemma~\ref{A:strict-normalization-rigidity}, applied to the
noetherian model $S$, gives
\[
 B\subset S^\nu\subset K.
\]
By Lemma~\ref{A:strict-overring-bound},
$\cd_p(\Spec S^\nu)\leq d$. Topological invariance of the
small \'etale site
\citep[Tag 04DZ, Theorem 59.45.2]{A_Stacks}
is an equivalence of sites for this integral universal
homeomorphism. It therefore compares all torsion coefficients
and gives
\[
 \cd_p(W)=\cd_p(\Spec S^\nu)\leq d.
\]
No replacement of an infinite normalization by a finite one has
been made in this comparison.

Proposition~\ref{A:spread} supplies a nonempty affine open on
which a nonzero generic degree-$d$ class spreads and remains
nonzero on every smaller nonempty affine open. Intersect with
$\Spec R[1/a]$ and with any prescribed nonempty open of
$\Spec R$, and choose a nonempty affine open $U$ in this
intersection. Every nonempty affine $U'\subset U$ then has
cohomological dimension at least $d$ by that generic class and at
most $d$ by the preceding paragraph. This proves all the
quantifiers of A.2. If $d=0$, the degree-zero generic class gives
the same conclusion.

For semilocal $R$, Lemma~\ref{A:semilocal-topnormalization}
provides a principal open whose normalization is a universal
homeomorphism. The preceding proof applies directly to that
open. Equivalently, the local-ring definition of geometric
unibranchness shows that this principal open has that property.
\end{proof}

\begin{remark}[Exact additional range]
This extends the normal-model result to geometrically unibranch
dense opens, and in particular to nonnormal semilocal models in
dimensions greater than two. It does not make the analogous
claim for every noetherian model. The finite-product argument
in Lemma~\ref{A:semilocal-topnormalization} uses semilocality:
an infinite product of local completions need not be noetherian,
so the cited EGA proposition cannot be applied to it. The
scope of these positive results is recorded in
Section~\ref{A:unresolved}. The unrestricted A.2 assertion is refuted
by the separate arithmetic construction below.
\end{remark}
